\documentclass[12pt]{article}

% setup configura os pacotes e os metadados
% --- Metadados do trabalho ---
\newcommand{\universidade}{Universidade Federal do Pará (UFPA)}
\newcommand{\instituto}{Instituto de Ciências Exatas e Naturais (ICEN)}
\newcommand{\faculdade}{Faculdade de Ciência da Computação}
\newcommand{\titulo}{Resenha: O Algoritmo de Babai para o Isomorfismo de Grafos}
\newcommand{\disciplina}{Grafos}
\newcommand{\professor}{Dr. Roberto Samarone dos Santos Araujo}

% Lista de autores separada por vírgula
\newcommand{\autores}{ Alessandro Reali Lopes Silva }

% E-mails separados por vírgula, dentro de chaves para formatação uniforme
\newcommand{\emails}{alessandro.silva@icen.ufpa.br}


% --- Idioma, codificacao e fonte base ---
\usepackage[utf8]{inputenc}
\usepackage[T1]{fontenc}
\usepackage[brazilian]{babel}
\usepackage{mathptmx}

% --- Matematica e simbolos ---
\usepackage{amsmath, amssymb, amsfonts}
\usepackage{siunitx}
\sisetup{group-separator = {.}, output-decimal-marker = {,}}
% --- Midia, graficos e codigo ---
\usepackage{graphicx}
\usepackage{xcolor}
\usepackage{listings}
\usepackage[font=small, labelfont=bf]{caption}
\usepackage{tikz}
\usepackage{pgfplots}
\pgfplotsset{compat=1.18}
\usepackage{booktabs}
\renewcommand{\arraystretch}{1.5}

% --- Layout e espacamento ---
\usepackage[a4paper, top=3cm, bottom=2cm, left=3cm, right=2cm]{geometry}
\usepackage[nopatch=footnote]{microtype}

\usepackage{titlesec}
\titlespacing*{\section}{0pt}{10pt}{10pt}
\titlespacing*{\subsection}{0pt}{8pt}{4pt}

\usepackage{parskip}

% --- Links e navegacao ---
\usepackage[colorlinks=true, linkcolor=black, urlcolor=blue]{hyperref}

% --- Estilo do Código (Listings) ---
\definecolor{codecomment}{rgb}{0.3, 0.3, 0.3} % Cinza escuro para comentários
\definecolor{codekeyword}{rgb}{0.0, 0.0, 0.0} % Preto negrito para keywords
\definecolor{codestring}{rgb}{0.2, 0.2, 0.2}  % Quase preto para strings

\lstset{
    basicstyle=\ttfamily\small,           % Fonte mono (Courier/Consolas)
    commentstyle=\itshape\color{codecomment},
    keywordstyle=\bfseries\color{codekeyword},
    stringstyle=\color{codestring},
    numbers=left,                         % Números à esquerda para referência
    numberstyle=\tiny\color{gray},
    stepnumber=1,
    numbersep=8pt,
    backgroundcolor=\color{white},
    showspaces=false,
    showstringspaces=false,
    showtabs=false,
    frame=lines,                          % Linha em cima e embaixo (estilo SBC)
    rulecolor=\color{black},
    tabsize=2,
    captionpos=b,                         % Legenda embaixo
    breaklines=true,                      % Quebra de linha automática
    breakatwhitespace=false,
    escapeinside={\%*}{*)},               % Para usar LaTeX dentro do código
    inputencoding=utf8,
    extendedchars=true,
    literate={á}{{\'a}}1 {é}{{\'e}}1 {í}{{\'i}}1 {ó}{{\'o}}1 {ú}{{\'u}}1 {ã}{{\~a}}1 {õ}{{\~o}}1 {ç}{{\c{c}}}1
}

% --- Macros e personalizacoes gerais ---
\newcommand{\divisoria}{\noindent\rule{\textwidth}{0.4pt}}
\newcommand{\sinal}[1]{\mathcal{S}(#1)}
\newcommand{\fourier}{\mathcal{F}}
\newcommand{\hz}{\text{Hz}}
\newcommand{\khz}{\text{kHz}}
\newcommand{\mhz}{\text{MHz}}
\newcommand{\ghz}{\text{GHz}}
\newcommand{\db}{\text{dB}}
\newcommand{\s}{\text{s}}
\newcommand{\ms}{\text{ms}}
\newcommand{\us}{\text{ }\mu\text{s}}
\newcommand{\bps}{\text{bps}}
\newcommand{\kbps}{\text{Kbps}}
\newcommand{\mbps}{\text{Mbps}}

\begin{document}

    \begin{center}

    {\Large \textbf{\titulo} \par}

    {\large \autores \par}

    {\small \texttt{\emails} \par}

    \divisoria

\end{center}


    O problema do Isomorfismo de Grafos (GI) representa um dos enigmas mais persistentes da ciência da computação,
    habitando uma zona de incerteza entre os problemas de classe P e os NP-Completos.
    O desafio central consiste em determinar se dois grafos com aparências distintas possuem, de fato, a mesma estrutura de conectividade subjacente.
    Durante mais de 30 anos, os algoritmos mais eficientes operavam em tempo subexponencial,
    um patamar que limitava significativamente a escalabilidade de soluções práticas para o problema.

    Em 2015, László Babai revolucionou a área ao propor um algoritmo que soluciona o GI em tempo quase-polinomial,
    aproximando-o consideravelmente da classe de problemas considerados "fáceis".
    O método emprega uma técnica de refinamento de cores (ou pintura de vértices), onde os nós são diferenciados com base em suas vizinhanças e propriedades de conectividade.
    O obstáculo técnico mais severo para essa abordagem são os Grafos de Johnson, cujas propriedades de supersimetria dificultam a identificação de discrepâncias,
    uma vez que todos os vértices parecem idênticos sob o esquema de pintura inicial.
    Babai demonstrou que, ao isolar e tratar especificamente esses casos, era possível contornar a simetria sem elevar o custo computacional a níveis proibitivos.

    O rigor da revisão científica foi evidenciado em 2017, quando o matemático Harald Helfgott identificou uma inconsistência no módulo "Split-or-Johnson" do algoritmo.
    O erro residia em uma recursão bifurcada que permitia o crescimento descontrolado de erros de aproximação, prejudicando a precisão do refinamento de cores em instâncias complexas.
    Ao substituir a bifurcação por uma chamada recursiva simplificada, Babai corrigiu a falha e validou o desempenho quase-polinomial,
    provando que o erro era uma falha de lógica modular e não um colapso do conceito fundamental.

    A transição do tempo subexponencial para o quase-polinomial representa um avanço estrutural na compreensão da complexidade algorítmica.
    Embora o GI ainda não tenha sido oficialmente inserido na classe P, o trabalho de Babai demonstra que problemas aparentemente intratáveis podem sucumbir a novas formas de redução de simetria.
    Essa descoberta não apenas impacta a teoria pura, mas também sinaliza uma evolução tecnológica capaz de transformar o processamento de redes complexas em diversas áreas da ciência e da engenharia.

    \divisoria
    
    \begingroup\raggedright
        \begin{thebibliography}{9}
            \bibitem{quanta2015}
            KLARREICH, Erica. \textbf{Landmark Algorithm Breaks 30-Year Impasse}. Quanta Magazine, 14 dez. 2015.
            Disponível em: \url{https://www.quantamagazine.org/algorithm-solves-graph-isomorphism-in-record-time-20151214/}.

            \bibitem{quanta2017}
            KLARREICH, Erica. \textbf{Graph Isomorphism Vanquished --- Again}. Quanta Magazine, 14 jan. 2017.
            Disponível em: \url{https://www.quantamagazine.org/graph-isomorphism-vanquished-again-20170114/}.
        \end{thebibliography}
    \endgroup

\end{document}
